\chapter{Mínimos Quadrados Ordinários}
\label{cap:ols}
% ── índice ──────────────────────────────────────────────────────
\index{mínimos quadrados ordinários|textbf}
\index{OLS|see{mínimos quadrados ordinários}}
\index{ols@\texttt{ols()}|textbf}
\index{viés de variável omitida|textbf}
\index{BLUE|textbf}
\index{consistência}
\index{heterocedasticidade}
\index{erros-padrão robustos}
\index{esttab@\texttt{esttab()}|textbf}

% ─────────────────────────────────────────────────────────────────────────────
\section{O modelo}

Considere o modelo de regressão linear:

\[
  y = X\beta + \varepsilon
\]

onde $y$ é um vetor $n \times 1$ de observações da variável dependente,
$X$ é uma matriz $n \times k$ de regressores (a primeira coluna é
identicamente 1, correspondendo ao intercepto), $\beta$ é um vetor
$k \times 1$ de parâmetros desconhecidos, e $\varepsilon$ é um vetor
$n \times 1$ de erros.

\subsection*{Hipóteses de Gauss-Markov}

\begin{enumerate}
  \item \textbf{Linearidade:} $y = X\beta + \varepsilon$.
  \item \textbf{Posto completo:} $\mathrm{rank}(X) = k$ (sem multicolinearidade perfeita).
  \item \textbf{Exogeneidade estrita:} $E[\varepsilon \mid X] = 0$.
  \item \textbf{Homocedasticidade:} $\mathrm{Var}[\varepsilon \mid X] = \sigma^2 I_n$.
  \item \textbf{Não-autocorrelação:} $\mathrm{Cov}[\varepsilon_i, \varepsilon_j \mid X] = 0$ para $i \neq j$.
\end{enumerate}

Sob H1–H3 o OLS é não-viesado. Sob H1–H5 é BLUE (Teorema de Gauss-Markov).

% ─────────────────────────────────────────────────────────────────────────────
\section{O estimador}

O estimador OLS minimiza a soma dos quadrados dos resíduos:

\[
  \hat{\beta} = \arg\min_{\beta}\, (y - X\beta)'(y - X\beta)
\]

A condição de primeira ordem $X'(y - X\beta) = 0$ — as \emph{equações normais} —
tem solução analítica:

\[
  \boxed{\hat{\beta} = (X'X)^{-1}X'y}
\]

Os resíduos são $\hat{e} = y - X\hat{\beta} = M_X y$, onde
$M_X = I - X(X'X)^{-1}X'$ é a matriz de resíduos (simétrica e idempotente).

% ─────────────────────────────────────────────────────────────────────────────
\section{Propriedades}

Sob H1–H3:
\[
  E[\hat{\beta}] = \beta
\]

Sob H1–H5, a variância condicional do estimador é:
\[
  \mathrm{Var}[\hat{\beta} \mid X] = \sigma^2 (X'X)^{-1}
\]

O estimador não-viesado da variância do erro é:
\[
  s^2 = \frac{\hat{e}'\hat{e}}{n - k}
\]

e portanto $\widehat{\mathrm{Var}}[\hat{\beta}] = s^2 (X'X)^{-1}$.

% ─────────────────────────────────────────────────────────────────────────────
\section{Verificação por DGP}

A prova mais direta de que o estimador está correto é definir um DGP
(\emph{data-generating process}) com parâmetros conhecidos, gerar dados
a partir dele e verificar que o OLS os recupera.

Escolha os parâmetros verdadeiros:
\[
  \beta_0 = 2{,}5 \qquad \beta_1 = 1{,}8 \qquad \beta_2 = -0{,}7
\]

e gere $n = 1000$ observações \emph{sem} termo de erro ($\varepsilon = 0$):
\[
  y_i = 2{,}5 + 1{,}8\,x_{1i} - 0{,}7\,x_{2i}
\]

Com $\varepsilon = 0$ as hipóteses de Gauss-Markov são satisfeitas de forma
exata, e o OLS deve recuperar $(\beta_0, \beta_1, \beta_2)$ com precisão
de máquina.

\begin{lstlisting}[caption={DGP sem ruído — verificação numérica}]
seed(42)
input df
id
1
end
for i in 1..=10 { df = append(df, df) }
df |> filter(_n <= 1000)
generate df x1 = rnormal()
generate df x2 = rnormal()
generate df y  = 2.5 + 1.8*x1 - 0.7*x2
ols(y ~ x1 + x2, df)
\end{lstlisting}

\noindent Saída real do interpretador:

\begin{lstlisting}[language={}, basicstyle=\ttfamily\scriptsize,
  frame=none, numbers=none, backgroundcolor=\color{codbg},
  xleftmargin=0pt, xrightmargin=0pt, breaklines=false]
=========================== OLS Regression Results ===========================
Dep. Variable:                     y || R-squared:                    1.0000
Model:                           OLS || Adj. R-squared:               1.0000
Covariance Type:          Non-Robust || F-statistic:                  0.0000
No. Observations:               1000 || Prob (F-statistic):         1.0000e0
Df Residuals:                    997 || Log-Likelihood:           31513.9419
AIC:                     -63021.8837 || BIC:                     -63007.1605

------------------------------------------------------------------------------
Variable   |       coef |    std err |        t |    P>|t| | [0.025  0.975]
------------------------------------------------------------------------------
const      |     2.5000 |     0.0000 | 6.17e+15 |    0.000 |  2.5000  2.5000
x1         |     1.8000 |     0.0000 | 3.36e+15 |    0.000 |  1.8000  1.8000
x2         |    -0.7000 |     0.0000 |-1.28e+15 |    0.000 | -0.7000 -0.7000
==============================================================================
\end{lstlisting}

O OLS recupera exatamente $(\hat\beta_0, \hat\beta_1, \hat\beta_2) = (2{,}5;\; 1{,}8;\; -0{,}7)$.
O erro padrão zero e o $R^2 = 1$ são esperados: sem ruído, os dados estão
exatamente sobre o hiperplano definido por $\beta$, e o estimador é numericamente
exato até a precisão de ponto flutuante.

\begin{nota}
  A estatística $F$ aparece como \texttt{0.0} e o p-valor como \texttt{1.0}
  neste caso — artefato numérico de $0/0$ quando $\mathrm{SQR} = 0$.
  As estatísticas $t$ são da ordem de $10^{15}$ pelo mesmo motivo: numerador
  finito dividido por denominador de precisão de máquina.
  Em presença de ruído ($\varepsilon \neq 0$) todas as estatísticas assumem
  valores finitos e interpretáveis.
\end{nota}

\subsection*{Efeito de má-especificação}

O mesmo DGP permite verificar o que acontece ao omitir $x_2$.
Estimam-se dois modelos e comparam-se com \lstinline{esttab}:

\begin{lstlisting}[caption={Modelo correto vs.\ $x_2$ omitido}]
let m1 = ols(y ~ x1 + x2, df)
let m2 = ols(y ~ x1, df)
esttab(m1, m2)
\end{lstlisting}

\begin{lstlisting}[language={}, basicstyle=\ttfamily\small,
  frame=none, numbers=none, backgroundcolor=\color{codbg},
  xleftmargin=0pt, xrightmargin=0pt, breaklines=false]
                            (1)              (2)
                            OLS              OLS
────────────────────────────────────────────────
x1                    1.8000***        1.7422***
                       (0.0000)         (0.0224)
x2                   -0.7000***
                       (0.0000)
Constant              2.5000***        2.1810***
                       (0.0000)         (0.0129)
────────────────────────────────────────────────
N                          1000             1000
R²                       1.0000           0.8585
Adj. R²                  1.0000           0.8583
────────────────────────────────────────────────
* p<0.10  ** p<0.05  *** p<0.01
\end{lstlisting}

\noindent O coeficiente de $x_1$ mal se altera (1.8000 $\to$ 1.7422) porque
$x_1$ e $x_2$ são gerados independentemente — a correlação amostral é
próxima de zero. Já a constante distorce de 2.5000 para 2.1810 porque
absorve $\hat\beta_2\,\bar{x}_2$, e a média amostral de $x_2$ não é
exatamente zero mesmo com $n=1000$.

\begin{nota}
  Viés de variável omitida (OVB) \emph{requer correlação} entre o regressor
  incluído e o omitido. Sem correlação, a omissão apenas infla os erros
  padrão e reduz o $R^2$ (perda de eficiência, não de consistência). Para o
  tratamento formal de endogeneidade e instrumentos, veja o
  Capítulo~\ref{cap:iv}.
\end{nota}

% ─────────────────────────────────────────────────────────────────────────────
\section{Erros padrão}

\subsection*{Clássico}

Assume homocedasticidade. A matriz de covariância estimada é
$s^2(X'X)^{-1}$ e o erro padrão do coeficiente $j$ é
$\mathrm{se}_j = s\,\sqrt{[(X'X)^{-1}]_{jj}}$.

\subsection*{Robusto à heterocedasticidade}

Quando $\mathrm{Var}[\varepsilon_i \mid x_i] = \sigma_i^2$ (heterocedasticidade),
o estimador sanduíche fornece inferência válida.

\textbf{HC1} (White com correção de graus de liberdade):
\[
  \hat{V}_{\mathrm{HC1}} = \frac{n}{n-k}\,(X'X)^{-1}
    \!\left(\sum_{i=1}^{n} \hat{e}_i^2\, x_i x_i'\right)\!(X'X)^{-1}
\]

\textbf{HC3} (mais conservador, recomendado em amostras pequenas):
\[
  \hat{V}_{\mathrm{HC3}} = (X'X)^{-1}
    \!\left(\sum_{i=1}^{n} \frac{\hat{e}_i^2}{(1 - h_{ii})^2}\, x_i x_i'\right)\!(X'X)^{-1}
\]

onde $h_{ii} = x_i'(X'X)^{-1}x_i$ é o \emph{leverage} da observação $i$.

\subsection*{Clusterizado}

Quando os erros são correlacionados dentro de grupos $g = 1,\ldots,G$:
\[
  \hat{V}_{\mathrm{cl}} = \frac{G(n-1)}{(G-1)(n-k)}\,
    (X'X)^{-1}\!\left(\sum_{g=1}^{G} X_g'\hat{e}_g\hat{e}_g'X_g\right)\!(X'X)^{-1}
\]

% ─────────────────────────────────────────────────────────────────────────────
\section{Inferência}

\subsection*{Teste $t$ sobre um coeficiente}

Sob H0: $\beta_j = \beta_j^0$:
\[
  t_j = \frac{\hat{\beta}_j - \beta_j^0}{\mathrm{se}(\hat{\beta}_j)}
  \;\xrightarrow{d}\; t_{n-k}
\]

\subsection*{Teste $F$ sobre restrições lineares}

Para $R\beta = r$ com $q$ restrições:
\[
  F = \frac{(R\hat{\beta} - r)'
       \bigl[R\,(X'X)^{-1}R'\bigr]^{-1}
       (R\hat{\beta} - r)}{q\,s^2}
  \;\xrightarrow{d}\; F_{q,\,n-k}
\]

\begin{nota}
  Com erros robustos ou clusterizados, $\hat{V}$ substitui $s^2(X'X)^{-1}$
  no cálculo do teste $F$.
\end{nota}

% ─────────────────────────────────────────────────────────────────────────────
\section{Sintaxe}

\begin{lstlisting}
ols(fórmula, df)
ols(fórmula, df, cov=tipo)
ols(fórmula, df, cov=cluster, cluster=coluna)
ols(fórmula, df, if = condição)
\end{lstlisting}

Alias: \lstinline{reg()} é equivalente a \lstinline{ols()}.

\section{Fórmula}
\label{sec:ols-formula}
\index{fórmula!sintaxe}
\index{fórmula!transformações}
\index{fórmula!interações}

A fórmula segue a notação \texttt{y \textasciitilde{} x1 + x2 + \ldots}:

\begin{lstlisting}[caption={OLS básico}]
input df
Y X
10 2
12 3
8  1
15 5
11 2
14 4
end

ols(Y ~ X, df)
\end{lstlisting}

A constante é incluída automaticamente. Para múltiplas variáveis:

\begin{lstlisting}
ols(Y ~ X1 + X2 + X3, df)
\end{lstlisting}

A fórmula pode ser passada como string:

\begin{lstlisting}
let formula = "Y ~ X1 + X2"
ols(formula, df)
\end{lstlisting}

\subsection{Transformações}
\label{subsec:formula-transformacoes}

Qualquer função suportada pela linguagem pode ser usada diretamente
no RHS da fórmula — o interpretador avalia a expressão
\emph{antes} de passar os dados ao estimador, de modo que o Greeners
recebe apenas colunas numéricas planas.

\begin{lstlisting}[caption={Transformações em regressores}]
// log de uma variável
ols(Y ~ log(X), df)

// raiz quadrada
ols(Y ~ sqrt(X), df)

// múltiplas transformações
ols(Y ~ log(K) + log(L), df)

// termo quadrático (equivalente a I(X^2))
ols(Y ~ X + X^2, df)
\end{lstlisting}

O nome do coeficiente na tabela de resultados reflete a expressão
original: \texttt{log(K)}, \texttt{sqrt(X)}, etc.

\subsection{Interações}
\label{subsec:formula-interacoes}

O operador \texttt{:} cria um termo de interação (produto
\emph{element-wise}) entre dois regressores. Cada lado pode ser uma
expressão arbitrária:

\begin{lstlisting}[caption={Interações}]
// interação entre variáveis simples
ols(Y ~ X1 + X2 + X1:X2, df)

// interação de transformações — log(K)*log(L)
ols(Y ~ log(K) + log(L) + log(K):log(L), df)

// apenas o termo de interação
ols(Y ~ log(K):log(L), df)
\end{lstlisting}

O operador \texttt{*} é açúcar sintático para expansão completa
(\texttt{x1*x2} $\equiv$ \texttt{x1 + x2 + x1:x2}):

\begin{lstlisting}
// equivalente: ols(Y ~ X1 + X2 + X1:X2, df)
ols(Y ~ X1 * X2, df)
\end{lstlisting}

\subsection{\texttt{I(expr)} — expressões aritméticas arbitrárias}
\label{subsec:formula-I}

A função \texttt{I(\ldots)} isola uma expressão aritmética composta
dentro da fórmula, protegendo os operadores \texttt{+}, \texttt{-} e
\texttt{*} do significado especial que têm na gramática de fórmulas:

\begin{lstlisting}[caption={Termos arbitrários com \texttt{I()}}]
// quadrado de X
ols(Y ~ X + I(X^2), df)

// interação linear construída explicitamente
ols(Y ~ X1 + X2 + I(X1 * X2), df)

// combinação linear de variáveis
ols(Y ~ I(X1 + X2), df)
\end{lstlisting}

\begin{nota}
  \texttt{I(X^2)} e \texttt{X\textasciicircum{}2} são equivalentes quando
  \texttt{X\textasciicircum{}2} aparece sozinho como termo (fora de
  interações e adições de nível superior). Use \texttt{I(\ldots)} para
  eliminar qualquer ambiguidade ou quando a expressão contém \texttt{+}
  ou \texttt{*}.
\end{nota}

\subsection{Variável categórica: \texttt{C()}}
\label{subsec:formula-categorica}

\texttt{C(var)} instrui o estimador a tratar \texttt{var} como
categórica, gerando dummies automáticas (categoria de referência = primeira
por ordem lexicográfica):

\begin{lstlisting}
ols(Y ~ X + C(regiao), df)
\end{lstlisting}

\section{Opções de erro padrão}

\begin{center}
\begin{tabular}{ll}
\toprule
\textbf{Opção} & \textbf{Fórmula} \\
\midrule
(padrão)                            & $s^2(X'X)^{-1}$ \\
\texttt{cov=robust}                 & HC1 (sanduíche White) \\
\texttt{cov=HC1}                    & idem \\
\texttt{cov=HC3}                    & HC3 (leverage-ajustado) \\
\texttt{cov=cluster, cluster=col}   & clusterizado por \texttt{col} \\
\bottomrule
\end{tabular}
\end{center}

\begin{lstlisting}[caption={Variantes de erro padrão}]
ols(Y ~ X, df)                           // clássico
ols(Y ~ X, df, cov=robust)              // HC1
ols(Y ~ X, df, cov=HC3)                 // HC3
ols(Y ~ X, df, cov=cluster, cluster=id) // clusterizado
\end{lstlisting}

\section{Amostra condicional}

A opção \lstinline{if =} restringe a estimação a um subconjunto sem modificar o DataFrame:

\begin{lstlisting}
ols(Y ~ X, df, if = grupo == 1)
\end{lstlisting}

\section{Capturando o resultado}

\begin{lstlisting}
let m = ols(Y ~ X, df)
display m
\end{lstlisting}

Quando atribuído, \lstinline{ols()} não imprime a tabela — apenas armazena o modelo.

\section{Pós-estimação}

\subsection{\texttt{predict}}

\begin{lstlisting}
let m = ols(Y ~ X, df)

// $\hat{y} = X\hat{\beta}$
let yhat = predict(m, "xb")
// $\hat{e} = y - X\hat{\beta}$
let ehat = predict(m, "residuals")
\end{lstlisting}

Aliases aceitos: \texttt{"fitted"} para valores ajustados; \texttt{"resid"} ou \texttt{"e"} para resíduos.

\subsection{\texttt{vif}}

O VIF do regressor $j$ é $\mathrm{VIF}_j = 1/(1 - R_j^2)$, onde $R_j^2$ é o
$R^2$ da regressão de $x_j$ sobre os demais regressores. Valores acima de 10
indicam multicolinearidade severa.

\begin{lstlisting}
let m = ols(Y ~ X1 + X2 + X3, df)
vif(m)
\end{lstlisting}

\subsection{\texttt{test}}

Teste $F$ de restrições lineares $H_0: R\beta = r$:

\begin{lstlisting}
let m = ols(Y ~ X1 + X2 + X3, df)
test(m, X2 = 0, X3 = 0)    // H0: X2 = X3 = 0
\end{lstlisting}

\subsection{\texttt{lincom}}

Combinação linear $c'\hat{\beta}$ com erro padrão $\sqrt{c'\hat{V}c}$:

\begin{lstlisting}
let m = ols(Y ~ X1 + X2, df)
lincom(m, X1 + X2)
\end{lstlisting}

\subsection{\texttt{margins}}

Efeitos marginais médios. No OLS linear $\partial E[y]/\partial x_j = \beta_j$,
equivalente aos coeficientes.

\begin{lstlisting}
let m = ols(Y ~ X, df)
margins(m)
\end{lstlisting}

\section{Tabela de resultados com \texttt{esttab}}

\begin{lstlisting}[caption={Comparação de especificações}]
let m1 = ols(Y ~ X, df)
let m2 = ols(Y ~ X, df, cov=robust)
let m3 = ols(Y ~ X1 + X2, df)

esttab(m1, m2, m3)
\end{lstlisting}

\begin{lstlisting}
eststo(ols(Y ~ X, df))
eststo(ols(Y ~ X1 + X2, df))
esttab()
\end{lstlisting}

\section{Bootstrap}

Erros padrão por bootstrap — alternativa não-paramétrica:

\begin{lstlisting}
bootstrap(ols, Y ~ X, df, n=1000)
bootse(m, n=500)
\end{lstlisting}
